\documentclass[10pt,a4paper]{amsart}
\usepackage[margin=19mm]{geometry}
\usepackage{amsmath,amssymb,mathtools}
\usepackage[scr=rsfs,cal=euler]{mathalfa}
\usepackage{xfrac}
\usepackage[unicode,hidelinks,bookmarks=false]{hyperref}
\newcommand{\ie}{i.e.\ }
\newcommand{\cf}{cf.\ }
\newcommand{\resp}{respectively\ }
\newcommand{\Subsection}{\subsection}
\providecommand{\nobreakdash}{\nobreak\hbox{-}}
\setlength{\emergencystretch}{2em}
\multlinegap0pt
\usepackage{mathrsfs}
\usepackage{color}
\usepackage{esint}
\allowdisplaybreaks
\usepackage{cleveref}
\newcommand{\on}{\operatorname}
\newcommand{\grad}{\on{grad}}
\newcommand{\divv}{\on{div}}
\newcommand{\tr}{\on{tr}}
\newcommand{\supp}{\on{supp}}
\newcommand{\cS}{\mathcal{S}}
\newcommand{\cH}{\mathcal{H}}
\newcommand{\spt}{\mathrm{spt}}
\newcommand{\reg}{\mathrm{reg}}
\newcommand{\sing}{\mathrm{sing}}
\newcommand{\hit}{\mathrm{hit}}
\newcommand{\merge}{\mathrm{merge}}
\DeclareMathOperator{\spread}{\mathrm{sp}}
\DeclareMathOperator{\strength}{\mathrm{str}}
\newcommand{\dee}{\partial}
\newcommand*{\ind}{\text{\usefont{U}{bbold}{m}{n}1}}
\DeclareMathOperator{\sign}{\mathrm{sgn}}
\newtheorem{lemma}{Lemma}
\newtheorem{proposition}[lemma]{Proposition}
\newtheorem{corollary}[lemma]{Corollary}
\newtheorem{theorem}[lemma]{Theorem}
\newtheorem{conjecture}[lemma]{Conjecture}
\newtheorem{remark}[lemma]{Remark}
\newtheorem{definition}[lemma]{Definition}
\newtheorem{example}[lemma]{Example}
\begin{document}
\title[A Partial epsilon-Regularity Theorem for Integral Stationary]{A Partial $\epsilon$-Regularity Theorem for Integral Stationary Geodesic Networks in $\mathbb{R}^n$}
\author{Henry Bosch}
\date{}
\maketitle
\noindent\emph{Source and licence.} A Partial $\epsilon$-Regularity Theorem for Integral Stationary Geodesic Networks in $\mathbb{R}^n$. Version arXiv:2607.23872v1 (CC BY 4.0). This material is distributed under the Creative Commons Attribution 4.0 International licence, http://creativecommons.org/licenses/by/4.0/, as recorded in the arXiv OAI licence metadata for this exact version and shown on the abstract page https://arxiv.org/abs/2607.23872v1. Full rights notice: licenses/arxiv-2607.23872-CC-BY-4.0.txt.\par\smallskip
\noindent\textit{Editorial scope.} Counted unit: the original \texttt{theorem} environment \texttt{structure-theorem} at source lines 162--171 together with its complete proof at lines 317--402; the delivered document keeps source lines 137--403 so that every referenced label and every citation resolves inside it. Native editable source is the paper's own concatenated TeX; the AMS wrapper is editorial formatting only. No publisher class, upstream script or unrelated artwork is redistributed.
\par\smallskip\noindent\textit{Reference note.} This article ships no compiled bibliography (biblatex); the entries delivered here are composed from the author's own BibTeX (.bib) fields, with no field invented and no entry dropped.
\begin{proposition}[First Variation] \label{first-variation}
Let $X \subset M$ be a stationary geodesic network with multiplicity. Let $U \subset\subset M$ be open, with the property that $X \cap U$ has finitely many ends, and each end has a unique limiting point in $\partial U$ ($U$ geodesically convex would suffice).
Let $Z$ be a $C^1$ vector field on $M$. Let $E$ be the set of ends of $X \cap U$. Each end $e \in E$ has a representative which is a geodesic segment $\gamma$ which limits to $\dee U$. For each $e \in E$ let $x_e \in \partial U$ be the limiting point in $\partial U$ of $\gamma$, let $\eta_e \in T_{x_e}M$ be the limiting unit tangent along $\gamma$ pointing outward from $U$, and let $\Theta(e) > 0$ be the multiplicity of $X$ along $\gamma$. Then
\[\int_U \divv_X Z \, dX = \sum_{e \in E} \Theta(e) \langle \eta_e, Z_{x_e}\rangle.\]
(Here for $x \in X$, $(\divv_X Z)|_x = \langle\eta, \nabla_\eta Z\rangle$, where $\eta \in T_xM$ is a choice of unit tangent to $X$ at $x$.)

\end{proposition}

\begin{definition}
    Let $X \subset M$ be a geodesic network with multiplicity. If $X$ is both integral and stationary we call it an \emph{integral stationary geodesic network}, or by abbreviation, an ISGN.
\end{definition}

\section{A Structure Theorem}
\label{sec:structuretheorem}

Fix $n \geq 2$. In this section, we prove a structure theorem for ISGNs sufficiently close to a multiplicity $k$ line in $\mathbb{R}^n$. In what follows, we let $E_1,\ldots,E_n$ be the standard orthonormal basis of $\mathbb{R}^n$ and of $T_x \mathbb{R}^n$ for any $x \in \mathbb{R}^n$. Let $x = (x_1,\ldots,x_n)$ be the standard coordinate on $\mathbb{R}^n$. Let $\pi = x_1$ be projection onto the $E_1$-axis.

\begin{definition}
    Let $B_r(0)$ denote the open ball of radius $r$ about the origin in $\mathbb{R}^{n-1}$. Suppose $A \subset \mathbb{R}$ is any subset. The \emph{tube} around $A$ of radius $r$ is the subset 
    \[T_{A,r} = A \times B_r(0) \subset \mathbb{R}^n.\]
    If $A$ is open, then $T_{A,r}$ is open. If $A = \{t\}$, write $T_{A,r} = T_{t,r}$.
\end{definition}

Now we can state the structure theorem for integral stationary geodesic networks in $\mathbb{R}^n$.


\begin{theorem}[Structure Theorem] \label{structure-theorem}
    Fix $k \geq 1$. For $i = 1,2,\ldots$ let $X_i$ be an ISGN in $\mathbb{R}^n$, and suppose that $X_i$ weakly converges to the multiplicity $k$ $E_1$-axis (as Radon measures on $\mathbb{R}^n$). Fix an open interval $I \subset\subset \mathbb{R}$ of finite length. Fix $0 < r < R$ and $\theta > 0$. Then there is an $N = N(n,k,(X_i), I, R, r, \theta)$ such that for all $i \geq N$,
    \begin{enumerate}
        \item $X_i \cap T_{I,R} \subset T_{I,r}$.
        \item At any regular point of $X_i \cap T_{I,R}$, any tangent vector to $X_i$ there makes an angle at most $\theta$ with the $E_1$-axis.
        \item At any singular point $x \in X_i \cap T_{I,R}$, the sum of the multiplicities of geodesic segments/rays entering $x$ from the left and leaving $x$ from the right (in the $E_1$ direction) are equal.
        \item For any $t \in I,$ $|X_i \cap T_{t,R}| \leq k$, and 
        \[\sum_{x \in X_i \cap T_{t,R}} \Theta_{X_i}(x) = k.\]
    \end{enumerate}
\end{theorem}


To prove the structure theorem, we need a series of lemmas. First, we prove that weak convergence implies local Hausdorff convergence.

\begin{lemma}[Local Hausdorff Convergence to the Line] \label{hausdorff-convergence}
    Suppose that $(X_i)$ is a sequence of ISGNs in $\mathbb{R}^n$ converging weakly to the multiplicity $k$ $E_1$-axis. Let $L \subset \mathbb{R}^n$ denote the $E_1$ axis, and let $K \subset \mathbb{R}^n$ be any compact set. Then
    \[\sup \, \{d(x,L): x \in K \cap X_i\} \to 0.\]
\end{lemma}

\begin{proof}
Assume not. After passing to a subsequence, there exists an $\epsilon > 0$ and $x_i \in X_i \cap K$ which satisfy $d(x_i,L) \geq 2\epsilon.$ By the monotonicity formula (see \cite{simon_gmt}, 4.3), we have $\mu_{X_i}(B_\epsilon(x_i)) \geq 2\epsilon$. Let $K_\epsilon = \{x \in \mathbb{R}^n: d(x,K) < \epsilon\}$. Then for any $x \in K$, $B_\epsilon(x) \subset K_\epsilon$.

Let $f \in C_c(\mathbb{R}^n)$ satisfy $f \geq 0$, $f(x) = 1$ if $x \in K_\epsilon$ and $d(x,L) \geq \epsilon$, and $f(x) = 0$ if $d(x,L) \leq \frac{\epsilon}{2}$. Then for any $i$,
\[\int_{\mathbb{R}^n} f \, dX_i \geq 2\epsilon,\]
but if $\mu$ is the measure associated to the multiplicity $k$-$E_1$ axis $L$,
\[\int_{\mathbb{R}^n} f \, d\mu = 0,\]
contradicting the weak convergence.
\end{proof}

Next we need a lemma that controls the integral tilt of an ISGN in terms of its distance to the $E_1$-axis.

\begin{definition}
    Let $\eta$ be a unit vector in $\mathbb{R}^n$. Its \emph{tilt} from the $E_1$-axis is the quantity $\sqrt{1 - \langle \eta,E_1\rangle^2}.$
\end{definition}

The tilt lies in $[0,1]$, is $0$ if $\eta$ is parallel to $E_1$ and is 1 if $\eta$ is perpendicular to $E_1$. Note that if $\phi = \angle(\eta, E_1)$ then the tilt is $\sqrt{1 - \langle\eta,E_1\rangle^2} = \sqrt{1 - \cos^2 \phi} = |\sin \phi|.$

\begin{lemma}[$L^2$ Tilt Control] \label{tilt-control}
Let $X$ be an ISGN in $\mathbb{R}^n$, and fix an open bounded $U \subset\subset \mathbb{R}^n$ which satisfies the hypotheses of Proposition \ref{first-variation}. As in that proposition, let $E$ be the set of ends of $X \cap U$, and for each $e \in E$, we have a point $x_e \in \partial U$, an outward pointing unit tangent $\eta_e \in T_{x_e}\mathbb{R}^n$, and a multiplicity $\Theta(e)$. For any $x \in \mathbb{R}^n$ let $r(x) = \sqrt{x_2^2 + \cdots + x_n^2}.$ Then
\[\int_U (1 - \langle \eta,E_1\rangle^2) \, dX \leq \sum_{e \in E} \Theta(e) \, r(x_e).\]
\end{lemma}

\begin{proof}
    Let $\eta$ be a choice of unit tangent on $X_\reg \cap U$ which is locally parallel on $X_\reg \cap U$. Note that $\eta$ may not agree with our choice of $\eta_e$ on the ends of $X$.

    Fix $j \in \{2,\ldots,n\}$ and consider the vector field $Z = x_jE_j$ on $\mathbb{R}^n$. Observe that on $X_\reg$, 
    \[\divv_X Z = \langle \eta, \nabla_\eta Z\rangle = \langle \eta, \eta[x_j] E_j + x_j \nabla_\eta E_j\rangle.\]
    But $\nabla_\eta E_j = 0$ and $\eta[x_j] = \langle \eta,E_j\rangle$ so 
    \[\divv_X Z = \langle \eta,E_j\rangle^2.\]
    Since $\eta$ is unique up to a sign, this is independent of our choice of $\eta$. Applying Proposition \ref{first-variation}, we obtain that 
    \[\int_U \langle \eta,E_j\rangle^2 \, dX = \sum_{e \in E} \Theta(e) (x_e)_j \langle \eta_e,E_j\rangle.\]
    Summing over $j$, and since $1 - \langle \eta,E_1\rangle^2 = \sum_{j=2}^n \langle \eta,E_j\rangle^2,$ we have 
    \[\int_U (1 - \langle \eta,E_1\rangle^2) \, dX = \sum_{e \in E} \Theta(e) \sum_{j=2}^n (x_e)_j \langle \eta_e,E_j\rangle \leq \sum_{e \in E} \Theta(e) \left|\sum_{j=2}^n (x_e)_j \langle \eta_e,E_j\rangle\right| \]
    \[\leq \sum_{e \in E} \Theta(e) \sqrt{\sum_{j=2}^n ((x_e)_j)^2\sum_{j=2}^n \langle \eta_e, E_j\rangle^2} \leq \sum_{e \in E} \Theta(e) r(x_e)\]
    and we are done.
\end{proof}

Next, let's guarantee that an ISGN which is contained in a tube and which only exits on the ends must span the whole tube.

\begin{lemma} \label{spans-tube}
Let $X$ be an ISGN in $\mathbb{R}^n$, $I \subset \subset \mathbb{R}$ an open interval of finite length, and $r > 0$. Suppose $X \cap T_{I,r} \neq \varnothing$ and $X \cap \partial T_{I,r} \subset T_{\partial I, r}$. Then, for each $t \in I$, $X \cap T_{t,r} \neq \varnothing$.  
\end{lemma}
\begin{proof}
    Suppose not. Since $I$ is connected, it suffices to show that the set $\pi(X \cap T_{I,r}) \subset I$ is both open and closed relative to $I$.
    
First, we show it is open. Fix $t \in \pi(X \cap T_{I,r})$, and let $Z$ be minimal $(X \cap T_{I,r})$-clopen set of $X \cap T_{I,r}$ containing $X \cap T_{t,r}$. If $Z \subset T_{t,r}$, then $Z$ must exit the tube $T_{I,r}$ along the neck at time $t$, contradiction. Thus there exists an $x \in Z$ at which there is a geodesic in $X$ leaving $x$ going positively or negatively in the $E_1$-direction. By the balancing condition, it follows that $X$ extends in both directions, which shows that $t$ lies in the interior of $\pi(X \cap T_{I,r})$.

Now we show it is closed. Fix $t_i \in \pi(X \cap T_{I,r})$ with $t_i \to t \in I$. Then select $x_i \in X \cap T_{t_i,r}$. Passing to a subsequence, $x_i \to x$, where $x \in \overline{T_{I,r}}$. But $\pi(x) = t \in I$ and $x \in X$. If $x \notin X \cap T_{I,r}$, then we must have $r(x)=r$, contradicting our assumption that $X \cap \partial T_{I,r} \subset T_{\partial I, r}$.
\end{proof}

Next, given an $L^2$ bound on the tilt, we want to find times which are regular, have small tilt, and have bounded total multiplicity. 

\begin{lemma} \label{find-point}
    Let $X$ be an ISGN in $\mathbb{R}^n$ and suppose that $I \subset\subset \mathbb{R}$ is open and bounded (not necessarily connected) with finite total length $L > 0$, and let $r > 0$. Suppose that for all $t \in I$, $X \cap T_{t,r}$ is nonempty. Further suppose that $\delta,C > 0$ and that we have an $L^2$ tilt bound
    \[\int_{T_{I,r}} (1 - \langle \eta,E_1\rangle^2) \, dX \leq \delta L\]
    as well as a mass bound
    \[\mu_X(T_{I,r}) \leq CL.\]
    Then there exists a $t^* \in I$ such that $X \cap T_{t^*,r}$ contains no singular points of $X$, and satisfies
    \[\sum_{x \in X \cap T_{t^*,r}} \Theta(x) \leq 3C\]
    and
    \[\max_{x \in X \cap T_{t^*,r}} \, (1 - \langle \eta_x, E_1\rangle^2) \leq 3\delta.\]
\end{lemma}

\begin{proof}
    For any $t \in I$ let
    \[m(t) = \sum_{x \in X \cap T_{t,r}} \Theta(x).\]
    By our assumption, we have $m(t) \geq 1$ for all $t \in I$. Note that if $t \in I$ has $ \#(X \cap T_{t,r}) = \infty$, then $m(t) = \infty$; this corresponds to segments of $X$ which are perpendicular to $E_1$. By local finiteness, $m(t) = \infty$ for only finitely many $t \in I$.
    
    For any $t \in I$, let
    \[A(t) = \sup_{x \in X \cap T_{t,r}} (1 - \langle \eta_x,E_1\rangle^2)\]
    if $X \cap T_{t,r}$ contains no singular points, and arbitrarily at other $t$ (of which there are finitely many).
    Fix $\epsilon > 0$ and let $J_1 = \{t \in I: m(t) > \frac{C}{\epsilon}\}.$ Since projection to the $E_1$ axis decreases length, we have
    \[\frac{C}{\epsilon}|J_1| \leq \int_{I} m(t) \, dt \leq \mu_X(T_{I,r}) \leq CL,\]
    so $|J_1| \leq \epsilon L$. Next let $J_2 = \{t \in I: A(t) > \frac{\delta}{\epsilon}\}.$ Similarly,
    \[\frac{\delta}{\epsilon}|J_2| \leq \int_{I} A(t) \, dt \leq \int_{T_{I,r}} (1 - \langle \eta,E_1\rangle^2) \, dX \leq \delta L,\]
    so $|J_2| \leq \epsilon L$. Taking $\epsilon = \frac{1}{3}$ we see that $|I \setminus J_1 \setminus J_2| \geq (1-\frac{2}{3})L > 0$. Since $X$ has finitely many singular points in $T_{I,r}$, take $t^* \in I \setminus J_1 \setminus J_2$ such that $X \cap T_{t^*,r}$ contains no singular points of $X$. Then $m(t^*) \leq 3C $ and $A(t^*) \leq 3\delta$, so we're done.
\end{proof}

The next lemma says that if we control the tilt on endpoints of a tube, then we control the tilt everywhere within the tube.

\begin{lemma} \label{brussel-sprout-lemma}
    Let $I = (t_1,t_2) \subset\subset \mathbb{R}$ be a bounded open interval, $r > 0$, $\delta > 0$, and $X \subset \mathbb{R}^n$ an ISGN which has the property that $X \cap \partial (T_{I,r}) \subset T_{\partial I, r}$. That is, $X$ exits the tube $T_{I,r}$ only ``along the ends''. As in Proposition \ref{first-variation} let $E$ be the set of ends of $X \cap T_{I,r}$. Suppose that for each $e \in E$, $r(\eta_e) \leq \delta$ (recall that $r(\eta) = \sqrt{\eta_2^2 + \cdots + \eta_n^2} = \sqrt{1 - \eta_1^2}$ is the tilt). Then at any regular point of $X \cap T_{I,r}$ any unit tangent $\eta$ to $X$ there has bounded tilt 
    \[r(\eta) \leq \delta \sum_{e \in E} \Theta(e).\]
\end{lemma}

\begin{proof}
Consider any $y \in X_\reg \cap T_{I,r}$, and let $\eta$ be a unit tangent to $X$ there at $y$. If $r(\eta) = 0$ there is nothing to prove, so assume that $r(\eta) > 0$. By a rotation of $\mathbb{R}^n$ which preserves $E_1$ (which leaves tilts invariant as well), we may assume that $\eta = (\eta_1,\eta_2,0,\ldots,0)$ where $\eta_2 > 0$. 

Let $R = \{x = (x_1,\ldots,x_n) \in T_{I,r}: x_2 < y_2\}$. Then $R$ is open and bounded in $\mathbb{R}^n$, so we can apply the first variation formula. Let $G$ be the set of ends of $X \cap R$. Then we can write $G = G_1 \sqcup G_2$, where $G_1 = \{e \in G: x_e \in T_{I,r}\}$ are ends whose points lie in the interior of $T_{I,r}$, and $G_2 = G \setminus G_1$. 
 Any $e \in G_1$ must exit $R$ along the face $\{x_2=y_2\} \cap T_{I,r} \subset \partial R$, so it must have $\langle \eta_e,E_2\rangle > 0$. Any $e \in G_2$ has $x_e \in \partial T_{I,r}$, so $x_e \in T_{\partial I,r}$ by our assumption, and we know that $r(\eta_e) \leq \delta$.
 
 Apply Proposition \ref{first-variation} with constant vector field $Z = E_2$.
We obtain that
\[0 = \sum_{e \in G} \Theta(e) \langle \eta_e,E_2\rangle = \sum_{e \in G_1} \Theta(e) \langle \eta_e ,E_2\rangle + \sum_{e \in G_2} \Theta(e) \langle \eta_e,E_2\rangle \geq \langle \eta_y,E_2 \rangle + \sum_{e \in G_2} \Theta(e)\langle \eta_e,E_2 \rangle. \]
Therefore
\[r(\eta_y) = \langle \eta_y,E_2\rangle \leq \sum_{e \in G_2} \Theta(e) | \langle \eta_e,E_2\rangle| \leq \delta \sum_{e \in E} \Theta(e) \]
so we are done.
\end{proof}

The next two lemmas show that if the angular deviation of $X$ is controlled, then the total multiplicity left of a singular point and right of a singular point must be equal. We choose to work with angles from now on instead of tilts, at the cost of a few constant factors.

\begin{lemma} \label{two-sides}
Let $U \subset\subset \mathbb{R}^n$ be open and let $X$ be an ISGN in $U$ with exactly one singular point $x$. Suppose further that no tangent to $X$ is perpendicular to $E_1$. Thus $X$ consists of a collection of geodesic segments $\mathcal{F}$ that enter $x$ from the left (in the $E_1$ direction) and a collection of geodesic segments $\mathcal{F}'$ that leave $x$ from the right. For a geodesic segment $I \in \mathcal{F} \cup \mathcal{F}'$, let $s_I \in S^{n-1}$ be the unique unit tangent vector to $I$ which points positively in the $E_1$ direction, let $\Theta(I)$ be the multiplicity of $I$, and let $M = \sum_{I \in \mathcal{F}} \Theta(I)$, $M' = \sum_{I' \in \mathcal{F}'} \Theta(I').$ If $M > M'$, then there exists $I \in \mathcal{F}$ such that $\angle(s_I,E_1) \geq \sqrt{\frac{2}{M}}$. Similarly if $M < M'$ then there exists $I' \in \mathcal{F}'$ such that $\angle(s_{I'},E_1) \geq \sqrt{\frac{2}{M'}}$.
    
\end{lemma}

\begin{proof}
   The stationarity condition asserts that
\[\sum_{I \in \mathcal{F}} \Theta(I) s_I = \sum_{I' \in \mathcal{F}'} \Theta(I') s_{I'}.\]
Without loss of generality assume that $M > M'$. Since $\langle s_{I},E_1\rangle \in (0,1]$ for all $I$, we have that
\[M' \geq \sum_{I' \in \mathcal{F}'} \Theta (I') \langle s_{I'} ,E_1\rangle = \sum_{I \in \mathcal{F}} \Theta(I)\langle s_{I},E_1\rangle.\]
Rearranging,
\[\sum_{I \in \mathcal{F}} \frac{\Theta(I)}{M}\langle s_{I},E_1\rangle \leq \frac{M'}{M}.\]
By properties of weighted averages, there exists $I \in \mathcal{F}$ such that
\[\langle s_I,E_1\rangle \leq \frac{M'}{M}.\]
Since $\cos \varphi \geq 1-\frac{\varphi^2}{2}$ for $\varphi \in [0,\frac{\pi}{2}]$, we get 
\[\frac{M'}{M} \geq \langle s_I,E_1\rangle \geq 1 - \frac{\angle(s_I,E_1)^2}{2},\]
so
\[\angle(s_I,E_1) \geq \sqrt{\frac{2(M-M')}{M}} \geq \sqrt{\frac{2}{M}}\]
as desired.
\end{proof}

\begin{lemma} \label{same-mult}
    Consider the same setup as in Lemma \ref{two-sides} but also assume that every unit tangent at a regular point of $X$ makes an angle at most $\theta$ with the $E_1$-axis. If $M \neq M'$, then there exists an $I \in \mathcal{F} \cup \mathcal{F}'$ such that 
    \[\angle(s_I,E_1) \geq \sqrt{\frac{2\cos \theta}{\min(M,M')}}.\]
\end{lemma}
\begin{proof}
    For any $I \in \mathcal{F} \cup \mathcal{F}'$, we have $\langle s_I,E_1\rangle \geq \cos \theta$. Without loss of generality, assume $M > M'$. Then we have
    \[M' \geq \sum_{I' \in \mathcal{F}'} \Theta(I') \langle s_{I'},E_1\rangle = \sum_{I \in \mathcal{F}} \Theta(I) \langle s_I,E_1\rangle \geq \cos \theta\sum_{I \in \mathcal{F}} \Theta(I) = M \cos \theta.\]

Thus $M \leq \frac{M'}{\cos \theta}$. From the conclusion of Lemma \ref{two-sides}, it follows that there exists $I \in \mathcal{F} \cup \mathcal{F}'$ such that 
\[\angle(s_I,E_1) \geq \sqrt{\frac{2}{M}} \geq \sqrt{\frac{2\cos \theta}{M'}}.\]
\end{proof}

Now we have enough to prove Theorem \ref{structure-theorem}.


\begin{proof}[Proof (of Theorem \ref{structure-theorem})] Fix $\epsilon \in (0,\frac{1}{4})$ and let $I' \subset \mathbb{R}$ be an $\epsilon$-neighborhood of $I$. That is, if $I = (a,b)$ then $I' = (a-\epsilon,b+\epsilon)$. Notice that given $r,\theta > 0$, we are free to choose $r', \theta'$ and prove the theorem for them, which implies the theorem for the given $r,\theta$. Thus we allow ourselves freely to decrease $r$ and $\theta$.

Decrease $\theta$ if necessary to ensure that
\[\theta < \sqrt{\frac{2\cos \theta}{3\frac{1+\epsilon}{1-4\epsilon}k}}\]
and that
\[ \cos \theta > \frac{k-1}{k}. \]
Then, decrease $r$ if necessary to ensure that 
\[12\sqrt{6} \left(\frac{1+\epsilon}{1-4\epsilon}\right)^{3/2}  \frac{k^{3/2}\sqrt{r}}{\epsilon} \leq \theta.\] 
Let $K = \overline{T_{I',R}}$. By local Hausdorff convergence of $X_i$ to the $E_1$-axis Lemma \ref{hausdorff-convergence} using $K$, we can choose $N$ such that for all $i \geq N$, $X_i \cap T_{I',R} \subset T_{I',r}$. 

The set $I' \setminus I$ has two connected components $(a-\epsilon,a]$ and $[b,b+\epsilon)$. Let $J_1 = (a-\epsilon,a)$ and $J_2 = (b,b+\epsilon)$ be their respective interiors. 

Then, by weak convergence of $X_i$ to the multiplicity $k$ $E_1$-axis, increase $N$ if necessary such that for all $i \geq N$ and $c \in \{1,2\}$,
\[(1-\epsilon)k |J_c| \leq \mu_{X_i}(T_{J_c,R}) \leq (1+\epsilon)k |J_c|.\]
In particular, this shows that $X_i \cap T_{I',r}$ is nonempty. Since $X_i \cap T_{I',R} \subset T_{I',r}$, $X_i$ cannot leave $T_{I',r}$ along its neck, so
by Lemma \ref{spans-tube}, for every $t \in I'$, $X_i \cap T_{t,r} \neq \varnothing$.

For any $t \in I'$, let $m_i(t) = \sum_{x \in X_i \cap T_{t,R}} \Theta(x)$. We make the following claim.

\ 

\textbf{Claim.} For all $i \geq N$, there exists an open interval $I''$ with $I \subset I'' \subset I'$ such that every tangent at a regular point of $X_i \cap T_{I'',R}$ makes an angle at most $\theta$ with the $E_1$-axis. Further, if $t^* \in \partial I''$ then
$m_i(t^*) \leq 3\frac{1+\epsilon}{1-4\epsilon} k.$

\

\emph{Sketch of proof of Claim.} Fix $i$. For $J_c$, $c \in \{1,2\}$, do the following. Use weak convergence to the multiplicity $k$ line to find two points $t_c < t_c'$ in $J_c$ close to either end of $J_c$ such that the local multiplicities $m_i(t_c)$ and $m_i(t'_c)$ are comparable to $k$. Then, use Lemma \ref{tilt-control} to show that the $L^2$ integral of tilt of $X_i$ over $(t_c,t'_c)$ is small. Then, plug that bound, along with a mass bound of $X_i$ over $(t_c,t'_c)$, into Lemma \ref{find-point} to obtain a point $t_c^*$ in $(t_c,t'_c) \subset J_c$ over which there is bounded multiplicity and bounded tilt. Let $I'' = (t_1^*,t_2^*).$ Since we have bounded tilt and multiplicity over $\partial I''$, we can apply Lemma \ref{brussel-sprout-lemma} to bound tilt everywhere on $I''$, which is equivalent to bounding the angle with the $E_1$-axis. 

\

\emph{Proof of Claim.} Fix $i$. Everything below will be carried out for $c \in \{1,2\}$. 
First we construct the points $t_c < t_c'$ in $J_c$. Observe that
\[(1+\epsilon)k|J_c| \geq \mu_{X_i}(T_{J_c,R}) \geq \int_{J_c} m_i(t) \, dt.\]
If we let
\[H_c = \left\{t \in J_c: m_i(t) > \frac{1+\epsilon}{\epsilon}k\right\},\]
then
\[(1+\epsilon)k|J_c| \geq \int_{J_c} m_i(t) \, dt \geq \frac{(1+\epsilon)k}{\epsilon}|H_c|,\]
so $|H_c| \leq \epsilon |J_c|.$ Thus $H_c$ cannot fill up either of the length $2\epsilon |J_c|$ intervals closest to either end of $J_c$. Let $t_c,t_c' \in J_c \setminus H_c$ be on the left and right sides of $J_c$, each within $2\epsilon |J_c|$ of $\partial J_c$. Then $t_c < t_c'$, $|t'_c - t_c| \geq (1-4\epsilon)|J_c| = \epsilon(1-4\epsilon)$, and $m_i(t_c),m_i(t'_c) \leq \frac{1+\epsilon}{\epsilon}k.$

Let $J'_c = (t_c,t_c')$. Next, we apply Lemma \ref{tilt-control} to the set $U_c = T_{J'_c,R}$, and we conclude that
\[\int_{U_c} (1-\langle \eta,E_1\rangle^2) \, dX_i \leq \sum_{e \in E_c} \Theta(e) r(x_e) \leq r\sum_{e \in E_c}\Theta(e) \leq 2kr \frac{1+\epsilon}{\epsilon},\]
where $E_c$ is the set of ends of $X_i \cap U_c$. Since $|J'_c| \geq \epsilon(1-4\epsilon),$ we have 
\[\int_{U_c} (1-\langle \eta,E_1\rangle^2) \, dX_i \leq 2kr\frac{1+\epsilon}{\epsilon} \leq 2kr \frac{1+\epsilon}{\epsilon^2(1-4\epsilon)} |J'_c|.\]

Now we can apply Lemma \ref{find-point} to the interval $J'_c$ with $\delta = \frac{2kr(1+\epsilon)}{\epsilon^2(1-4\epsilon)}$. Since
\[\mu_{X_i}(T_{J'_c,r}) \leq \mu_{X_i}(T_{J_c,r}) \leq (1+\epsilon)k |J_c| = \epsilon(1+\epsilon)k \leq \frac{1+\epsilon}{1-4\epsilon}k |J'_c|,\]
we can take $C = \frac{1+\epsilon}{1-4\epsilon}k$ in the mass bound. Therefore, by Lemma \ref{find-point}, there exists a $t^*_c \in J'_c \subset J_c$ such that the sets $X_i \cap T_{t_c^*,r}$ for $c=1,2$ have no singular points of $X_i$ and satisfy

\[\sum_{x \in X_i \cap T_{t_c^*,r}
} \Theta(x) \leq 3\frac{1+\epsilon}{1-4\epsilon} k\]
and
\[\max_{x \in X_i \cap T_{t_c^*,r}} (1-\langle \eta_x,E_1\rangle^2) \leq \frac{6(1+\epsilon)}{(1-4\epsilon)}\frac{kr}{\epsilon^2}.\]

Let $I'' = (t_1^*,t_2^*)$. Apply Lemma \ref{brussel-sprout-lemma}, where \[ \delta = \sqrt{\frac{6(1+\epsilon)}{(1-4\epsilon)}}\frac{\sqrt{kr}}{\epsilon} \]
is the upper bound on tilt on the endpoints. We conclude that every tilt at a regular point of $X_i \cap T_{I'',R}$ is at most

\[6\sqrt{6} \left(\frac{1+\epsilon}{1-4\epsilon}\right)^{3/2}  \frac{k^{3/2}\sqrt{r}}{\epsilon}. \]

Since $|\phi| \leq 2|\sin \phi|$ for any $\phi \in [0,\frac{\pi}{2}]$, we can bound the angle by two times the tilt. Therefore at every regular point of $X_i \cap T_{I'',R}$, the angle against the $E_1$-axis is at most

\[12\sqrt{6} \left(\frac{1+\epsilon}{1-4\epsilon}\right)^{3/2}  \frac{k^{3/2}\sqrt{r}}{\epsilon} \]
which is at most $\theta$ by assumption. That completes the proof of the claim.

\

Let's continue with the proof of Theorem \ref{structure-theorem}. To prove (3), suppose for contradiction it is false, and let $x \in X_i \cap T_{I'',R}$ be a singular point failing the condition in (3) which is minimal with respect to the $x_1$-coordinate (such a point exists, since by local finiteness, $X_i$ has finitely many singular points in $T_{I'',R}$). By minimality, the total multiplicity $M$ coming into $x$ from the left is at most $m_i(t_1^*) \leq 3\frac{1+\epsilon}{1-4\epsilon}k$. This contradicts Lemma \ref{same-mult}, by our assumption that $\theta < \sqrt{\frac{2\cos \theta}{3\frac{1+\epsilon}{1-4\epsilon}k}}$.

\ 

Finally we move to prove (4). From the previous point, we know that $m_i(t)$ is constant on $I''$, so it is constant on $I$. Since all tangents are within an angle $\theta$ with the $E_1$-axis, it follows that
\[m_i |I| \leq \mu_{X_i}(T_{I,R}) \leq \frac{m_i|I|}{\cos \theta}.\]
Letting $a_i = \frac{1}{|I|}\mu_{X_i}(T_{I,R}),$
we have that
\[m_i \leq a_i \leq \frac{m_i}{\cos \theta}.\]
 By weak convergence we have
\[a_i|I| = \mu_{X_i}(T_{I,R}) \to k|I|,\]
so $a_i \to k$. 
Therefore 
\[\liminf_{i \to \infty}m_i \geq \liminf_{i \to \infty} a_i \cos \theta  = k \cos \theta \]
and
\[\limsup_{i \to \infty}m_i \leq \limsup_{i \to \infty} a_i = k.\]
But the $m_i$ are integers. Since we assumed that $\cos \theta > \frac{k-1}{k}$, we get that 
$\liminf_{i \to \infty} m_i > k-1,$
so
$\liminf_{i \to \infty} m_i = k,$ so $\lim_{i \to \infty}m_i = k.$ Increase $N$ again if necessary such that for all $i \geq N$, $m_i = k$, and we are done.
\end{proof}


\begin{thebibliography}{99}
\bibitem[]{simon_gmt} Simon, Leon. Lectures on Geometric Measure Theory. 3 (1984)
\end{thebibliography}
\end{document}
